% =====================================================================
%  DLoc Lightweight Localization — COMPLETE PROJECT GUIDE
%  The one reference: recreate the full environment & folder, the GitHub
%  workflow, running experiments, and the NILES paper + slides.
%  Compile with Tectonic ( tectonic Team_Setup_Guide.tex ) or pdfLaTeX.
% =====================================================================
\documentclass[11pt]{article}

\usepackage[a4paper,margin=0.95in]{geometry}
\usepackage{xcolor}
\usepackage{listings}
\usepackage{enumitem}
\usepackage{booktabs}
\usepackage{parskip}
\usepackage[hidelinks]{hyperref}

\definecolor{cmdbg}{gray}{0.95}
\lstdefinestyle{sh}{
  basicstyle=\ttfamily\footnotesize, backgroundcolor=\color{cmdbg},
  frame=single, framerule=0pt, rulecolor=\color{cmdbg},
  breaklines=true, columns=fullflexible, keepspaces=true,
  xleftmargin=6pt, xrightmargin=2pt, aboveskip=7pt, belowskip=7pt,
  showstringspaces=false, literate={~}{{\textasciitilde}}1,
}
\lstset{style=sh}
\setlength{\emergencystretch}{3em}\sloppy
\newcommand{\code}[1]{\lstinline[basicstyle=\ttfamily]|#1|}

\title{\textbf{DLoc Lightweight WiFi Localization}\\[2pt]
\large Complete Project Guide --- environment, GitHub, experiments, and the NILES paper/slides}
\author{E-JUST graduation project}
\date{\today}

\begin{document}
\maketitle
\vspace{-1.1em}\hrule\vspace{1em}

This is the single reference for the whole project. By the end you will have an
\textbf{exact copy of the working folder} --- the code, the $\sim$70\,GB WILD dataset,
the trained model weights, the shared Claude ``memory'', and all of the paper/thesis/
slide sources for NILES. Everything you need to present and to continue the work.

Follow the steps \textbf{in order}. You already have all the access (GitHub write +
the HuggingFace org); you only have to pull everything down. Budget $\sim$30--60\,min,
mostly the dataset download.

% =====================================================================
\section{Where everything lives}
% =====================================================================
The project is split across a code repo (small, in git) and HuggingFace (the big
files). A \texttt{git clone} gives you the code and \emph{all the writing/slides};
you pull the dataset and weights from HuggingFace.
\begin{itemize}[leftmargin=1.4em,itemsep=3pt]
\item \textbf{GitHub} --- \texttt{Abdulrahman-bits2/DLoc-lightweight-localization}:
      code, paper, thesis, slides, figures, reports. \emph{You have write access.}
\item \textbf{HF} \texttt{JustAGeek/dloc-wild-fig10b}: the WILD dataset ($\sim$70\,GB). Public.
\item \textbf{HF} \texttt{JustAGeek/dloc-code}: trained weights, logs, results. Public.
\item \textbf{HF} \texttt{Graduation-project123/dloc-private}: the Claude project
      memory. Private --- reached through your org membership (\texttt{Ahmed-Emad-Haggag}).
\end{itemize}

% =====================================================================
\section{Prerequisites (install once)}
% =====================================================================
\begin{itemize}[nosep,leftmargin=1.4em]
\item \textbf{git}, \textbf{Python 3.9+} (3.10--3.13 fine) with \texttt{pip}.
\item Your \textbf{HuggingFace account} \texttt{Ahmed-Emad-Haggag} (already in the org).
\item A \textbf{GPU + CUDA} \emph{only to train}. To read code, reproduce the published
      numbers from the downloaded weights, build the paper, or present --- no GPU needed.
\item To \textbf{edit/compile the paper} locally: a LaTeX install (TeX Live / MiKTeX)
      or just use \href{https://overleaf.com}{Overleaf} (Section~8). To open the
      \textbf{slides}: Microsoft PowerPoint.
\end{itemize}

% =====================================================================
\section{Get the full environment \& folder --- the easy way}
% =====================================================================
Four steps. The repo ships \texttt{setup.sh} that does the rest.

\textbf{Step 1 --- log in to HuggingFace} (needed for the private memory repo):
\begin{lstlisting}
pip install -U "huggingface_hub[cli]"
hf auth login        # paste a token from huggingface.co/settings/tokens (Read is enough)
hf auth whoami       # should print: Ahmed-Emad-Haggag
\end{lstlisting}

\textbf{Step 2 --- clone the code:}
\begin{lstlisting}
git clone https://github.com/Abdulrahman-bits2/DLoc-lightweight-localization.git
cd DLoc-lightweight-localization
\end{lstlisting}

\textbf{Step 3 --- run the setup script:}
\begin{lstlisting}
bash setup.sh                 # asks before the ~70 GB dataset download
#   bash setup.sh --all -y    # everything, no questions (includes the dataset)
#   bash setup.sh --no-dataset # skip the dataset (code + weights + memory only)
#   bash setup.sh --with-teacher  # also regenerate teacher outputs (needs a GPU)
\end{lstlisting}
It makes a virtual environment, installs dependencies, and downloads the dataset,
weights, and memory, then creates the cross-session symlinks.

\textbf{Step 4 --- check it worked:}
\begin{lstlisting}
ls code/data code/weights claude-memory
\end{lstlisting}
You should see the \texttt{.mat} dataset files, 10 \texttt{.pth} weights, and 20 memory
\texttt{.md} files. That is the full folder --- same as the original.

% =====================================================================
\section{Get the environment --- the manual way (alternative)}
% =====================================================================
Same result, by hand, from the cloned folder:
\begin{lstlisting}
python -m venv .venv
source .venv/bin/activate              # Windows: .venv\Scripts\activate
pip install -r code/requirements.txt

cd code
hf download JustAGeek/dloc-wild-fig10b --repo-type dataset --local-dir data   # ~70 GB
hf download JustAGeek/dloc-code --include "weights/*" --local-dir .           # weights/
cd ..
hf download Graduation-project123/dloc-private --repo-type dataset --local-dir .  # claude-memory/

# cross-session symlinks (only needed for the LOSO runs):
cd code/data
ln -s dataset_jacobs_Aug16_1.mat     dataset_train_jacobs_Aug16_1.mat
ln -s dataset_jacobs_Aug16_3.mat     dataset_train_jacobs_Aug16_3.mat
ln -s dataset_jacobs_Aug16_4_ref.mat dataset_train_jacobs_Aug16_4_ref.mat
cd ../..
\end{lstlisting}
The Mamba model only, if you ever run it, needs two extra CUDA-only packages:
\code{pip install mamba-ssm --no-deps} then \code{pip install causal-conv1d}.

% =====================================================================
\section{GitHub workflow (how we collaborate on the code)}
% =====================================================================
When you cloned from GitHub, your \texttt{origin} remote \emph{is} the GitHub repo, so
the normal workflow just works. \textbf{Always pull before you start}, and \textbf{work
on a branch}, not directly on \texttt{main}.
\begin{lstlisting}
git pull                                   # get everyone's latest changes first

git checkout -b ahmed/my-change            # start your own branch
#   ... edit files ...
git add -A
git commit -m "short description of what I changed"
git push -u origin ahmed/my-change         # pushes to GitHub (origin = GitHub)
# then open a Pull Request on GitHub to merge your branch into main
\end{lstlisting}
Useful:
\begin{lstlisting}
git status                 # what changed
git log --oneline -10      # recent history
git checkout main && git pull   # back to main and update
\end{lstlisting}
\textbf{Rules:}
\begin{itemize}[nosep,leftmargin=1.4em]
\item \textbf{Treat HuggingFace as download-only} --- use \code{hf download}, never
      \texttt{git push} to a HuggingFace remote. The HF repos hold the only copies of
      the weights/dataset/memory; a force-push can destroy them.
\item \textbf{Don't commit big files} (datasets, \texttt{.pth}/\texttt{.pt}, logs, run
      folders). They are git-ignored on purpose --- they live on HuggingFace.
\end{itemize}

% =====================================================================
\section{The project memory + your own Claude}
% =====================================================================
\texttt{claude-memory/} is the project's full history --- every run, data split, result,
decision, and the paper/slide status (20 notes, indexed by \texttt{MEMORY.md}). To pick
up exactly where things are:
\begin{enumerate}[nosep,leftmargin=1.5em]
\item Open \textbf{Claude Code} (your own account) \emph{inside} the project folder.
\item Tell it: \emph{``Read \texttt{claude-memory/MEMORY.md} and the files it indexes ---
      that is this project's full context.''}
\item Claude now knows the whole project (including all the NILES/paper decisions).
\end{enumerate}

% =====================================================================
\section{Running experiments}
% =====================================================================
Run everything from the \texttt{code/} folder with the venv active.
\begin{lstlisting}
cd code
source ../.venv/bin/activate              # Windows: ..\.venv\Scripts\activate

# --- verify the PUBLISHED results WITHOUT a GPU (uses the downloaded weights) ---
python evaluation/benchmark_models.py      # efficiency table (needs no data)
python evaluation/test_robustness.py       # set MODEL_TYPE + WEIGHTS_PATH at top of file
python evaluation/eval_cross_session.py     # set MODEL_TYPE + WEIGHTS_PATH at top of file

# --- retrain from scratch (needs a GPU) ---
python training/precompute_teacher.py       # once, before any DISTILLED run
python training/train_distillation.py        # students: edit STUDENT_TYPE / MODE at top
python training/train_unet_advanced.py        # U-Net:  edit SCHEDULER_TYPE / MODE at top
bash    training/run_crosssession_table1.sh    # LOSO: 3 folds x 3 seeds
\end{lstlisting}
The weight files to use for each model are listed in \texttt{REPRODUCIBILITY.md}
(the \texttt{MODEL\_TYPE}/\texttt{WEIGHTS\_PATH} table). Every data split and every run
is documented in \texttt{Experiments\_and\_Data\_Splits.pdf}.

% =====================================================================
\section{The NILES paper \& slides (for your presentation)}
% =====================================================================
Everything for the talk is already in the repo (you got it with \texttt{git clone}).

\textbf{The slides.} Open in PowerPoint:
\begin{lstlisting}
NilesConf/WiFi_Localization_NILES_v6.pptx
\end{lstlisting}
This is the final 38-slide deck (backup slides merged in, all numbers checked).

\textbf{The paper (camera-ready).} The compiled PDF is ready to read:
\begin{lstlisting}
NilesConf/camera_ready_latex/main.pdf       # the camera-ready paper (read this)
NilesConf/camera_ready_latex/main.tex       # its LaTeX source (+ references.bib, figures/)
NilesConf/submission/main.tex               # same paper, the submission source
\end{lstlisting}
To \textbf{edit \& recompile} the paper, two options:
\begin{itemize}[nosep,leftmargin=1.4em]
\item \textbf{Overleaf (easiest):} new project $\rightarrow$ \emph{Upload} all files from
      \texttt{NilesConf/camera\_ready\_latex/} $\rightarrow$ compile. It is self-contained
      (includes \texttt{IEEEtran.cls}, \texttt{references.bib}, \texttt{figures/}).
\item \textbf{Locally} (from inside that folder):
\begin{lstlisting}
cd NilesConf/camera_ready_latex
# with Tectonic (one command, handles the bibliography):
tectonic main.tex
# or with a standard TeX Live / MiKTeX install:
pdflatex main && bibtex main && pdflatex main && pdflatex main
\end{lstlisting}
\end{itemize}

\textbf{Figures \& numbers} (useful while presenting):
\begin{itemize}[nosep,leftmargin=1.4em]
\item \texttt{Figures/generated/} --- all publication figures (CDFs, heatmaps, compression).
\item \texttt{Experiments\_and\_Data\_Splits.pdf} --- every data split and H100 run, precisely.
\item \texttt{reports/phase2\_results.tex} + \texttt{reports/CROSSSESSION\_FINAL\_REPORT.md}
      --- the canonical result tables (medians/P90/P99, LOSO vs DLoc).
\end{itemize}

\textbf{Also in the repo:} \texttt{NilesConf/journal\_paper/} (extended journal version,
\texttt{Journal\_paper.pdf}), \texttt{thesis/} (full thesis + \texttt{Final\_Thesis.pdf}),
\texttt{presentation/} (an HTML defense-deck variant), \texttt{guide/} (study guide).

% =====================================================================
\section{What the folder looks like (reference)}
% =====================================================================
After setup, the folder matches the original:
\begin{lstlisting}
DLoc-lightweight-localization/
  setup.sh                         one-shot environment setup
  ONBOARDING.md  REPRODUCIBILITY.md    guides (result -> command map)
  Team_Setup_Guide.pdf             <- this guide
  Experiments_and_Data_Splits.pdf  every data split + H100 run
  code/
    core/ models/ training/ evaluation/ configs/ figures/   source
    data/           WILD dataset .mat files        (from HF, after setup)
    weights/        10 trained .pth checkpoints     (from HF, after setup)
    results_crosssession/   final LOSO result JSONs
    requirements.txt
  claude-memory/                   the Claude project memory (from HF, private)
  data_split_idx/                  the 10 original split-index .mat files
  Figures/        generated/ (figures)  figure_data/ (per-sample arrays)
  reports/                         LaTeX/Markdown result reports
  NilesConf/
    WiFi_Localization_NILES_v6.pptx   <- the NILES slides
    camera_ready_latex/   main.tex  main.pdf  references.bib  figures/   <- the paper
    submission/           main.tex (camera-ready source)
    journal_paper/        extended journal version (Journal_paper.pdf)
    NILES2026-CFP.pdf     the call-for-papers
  thesis/  paper/  guide/  presentation/     writing artifacts
  .venv/                           Python environment (after setup)
\end{lstlisting}

% =====================================================================
\section{Troubleshooting}
% =====================================================================
\begin{description}[leftmargin=0pt,style=unboxed,itemsep=3pt]
\item[\texttt{hf: command not found}] \code{pip install -U "huggingface_hub[cli]"} and
      make sure the venv is active.
\item[Private memory ``403''/``not found''] \code{hf auth whoami} must print
      \texttt{Ahmed-Emad-Haggag}, and you must be in the \texttt{Graduation-project123}
      org (you are).
\item[Dataset download interrupted] Re-run it --- \code{hf download} resumes. Skip it
      with \code{bash setup.sh --no-dataset}.
\item[LOSO error about \texttt{dataset\_train\_jacobs\_Aug16\_*}] The symlinks are
      missing (Section~5); \texttt{setup.sh} makes them. On Windows, enable Developer
      Mode or use \texttt{mklink}.
\item[\texttt{git push} goes to the wrong place] From a GitHub clone, \texttt{origin}
      is GitHub --- that is correct. Never add or push to a HuggingFace git remote.
\item[Windows] Activate the venv with
      \texttt{.venv\textbackslash Scripts\textbackslash activate}; symlinks need
      Developer Mode. The shared GPU workstation is Linux, where all commands work as written.
\end{description}

% =====================================================================
\section*{Appendix --- command cheat-sheet}
% =====================================================================
\addcontentsline{toc}{section}{Appendix --- command cheat-sheet}
\begin{lstlisting}
# ---- one-time full setup -------------------------------------------
pip install -U "huggingface_hub[cli]" && hf auth login
git clone https://github.com/Abdulrahman-bits2/DLoc-lightweight-localization.git
cd DLoc-lightweight-localization && bash setup.sh --all -y

# ---- daily git ------------------------------------------------------
git pull
git checkout -b ahmed/my-change
git add -A && git commit -m "..." && git push -u origin ahmed/my-change

# ---- run (from code/, venv active) ---------------------------------
cd code && source ../.venv/bin/activate
python evaluation/benchmark_models.py
python evaluation/eval_cross_session.py        # set MODEL_TYPE + WEIGHTS_PATH

# ---- build the paper ------------------------------------------------
cd NilesConf/camera_ready_latex && tectonic main.tex     # or upload the folder to Overleaf

# ---- open the slides ------------------------------------------------
# NilesConf/WiFi_Localization_NILES_v6.pptx   (PowerPoint)
\end{lstlisting}

\vspace{0.6em}\hrule\vspace{0.5em}
Questions? Ping the team, or open the project in Claude Code and point it at
\texttt{claude-memory/MEMORY.md}. Good luck with the NILES talk!

\end{document}
